\section{Finite-level arithmetic Selmer estimates and auxiliary limits}
\label{sec:sr-arithmetic-selmer-producers}
The statements below construct the finite trace detectors and the
auxiliary arithmetic tower, and prove the first- and second-Selmer
bounds needed for the common-support argument.  They apply to the
actual nice paired points with their displayed local conditions.
Their constants may depend on that fixed finite family of points;
the uniform conductor estimates later absorb this dependence through
the common support and automorphic counting argument.


\subsection{The arithmetic Taylor--Wiles tower and its common limit}
\label{sec:arithmetic-tw-tower}

We construct the actual auxiliary arithmetic groups and the limit
used by the completed-support argument.
The number of auxiliary primes is fixed by the selected trace
detectors; it need not equal a conjectural dual Selmer dimension.

Let \(L/L^+\) be the CM extension in the unitary problem, let
\(\ell\) be odd, and let \(S\) contain all places above \(\ell\)
and all ramification of the fixed residual packet.
Write \(\Gamma_0\) for the arithmetic quotient over \(L^+\),
\(G_0=\ker\eta\) for its CM subgroup, and
\(G_0'=\ker(\varepsilon|_{G_0})\), using the \emph{full}
\(\ell\)-adic cyclotomic character.
Keep the actual residual characters \(\bar\chi_1,\ldots,\bar\chi_n\)
and their polarized determinant.
Suppose a finite set
\[
                   s_1,\ldots,s_r\in G_0'
 \tag{TW1}
\]
has been chosen.
For each of the finitely many nice paired points under
consideration, require the characteristic polynomials
of \(\rho(s_j)\) to be separable over its fraction field.
The trace-detection proof chooses the \(s_j\) once, before the
tower.  Regularity is used in the Selmer estimates, not in
Chebotarev's construction below.
The source of the regular-coset property for the actual curve
points must therefore be established at the nice-point step;
absolute irreducibility on the full cyclotomic kernel alone is
not substituted for it.

We use only the following standard arithmetic finiteness input:
for every open subgroup \(U\) of the global arithmetic group,
and of each fixed local group, the set
\(\operatorname{Hom}_{\mathrm{cts}}(U,\mathbb F_\ell)\) is finite.
Globally, an open subgroup corresponds to a finite number-field
extension with a fixed finite ramification set.
Kummer theory, finite generation of its \(S\)-unit group, and
finiteness of its class group give this assertion after adjoining
\(\mu_\ell\); restriction and corestriction, or the
inflation--restriction sequence, give it before adjoining
\(\mu_\ell\).
Locally it follows from local class field theory.
This is the actual \srcid{PhiCond} input; it is independent of
the coefficient field and of the auxiliary tower level.

\paragraph{A residual-kernel subgroup and a cofinal sequence.}
Set
\[
 H=\{g\in\Gamma_0:\eta(g)=1,\ 
          \bar\chi_i(g)=1\ (1\leq i\leq n),\
          \bar\varepsilon(g)=1\}.
 \tag{TW2}
\]
It is open and normal.
Normality under \(G_0\) follows because the \(\bar\chi_i\) are
characters.  Under complex conjugation, polarization identifies
the multiset of conjugate characters with their inverses times
the cyclotomic multiplier; on the displayed kernel the multiplier
is \(1\), so that kernel is preserved.

For an open normal \(J\subset\Gamma_0\), put
\[
 \Phi_\ell(J)=
       \bigcap_{\alpha\in
          \operatorname{Hom}_{\mathrm{cts}}(J,\mathbb F_\ell)}
                         \ker\alpha .
\]
The intersection is finite, so \(\Phi_\ell(J)\) is open;
it is characteristic in \(J\) and hence normal in \(\Gamma_0\).
Take \(W_0=H\) and \(W_{N+1}=\Phi_\ell(W_N)\).
Define, exactly as in the source,
\[
 \mathrm{PL}_\ell(H)=
 \bigcap_{\substack{V\triangleleft_{\mathrm{open}}\Gamma_0\\
     \forall h\in H\ \exists e\geq0:\ h^{\ell^e}\in V}}V.
 \tag{TW3}
\]
Thus a finite quotient used in this intersection has an
\(\ell\)-group image of \(H\).
We have \(W_N\subset H\), and
\[
 \forall V\triangleleft_{\mathrm{open}}\Gamma_0,\
 \mathrm{PL}_\ell(H)\subset V
 \quad\Longrightarrow\quad
 W_N\subset V\quad(N\text{ sufficiently large}).
 \tag{TW4}
\]
To check this, the image \(P\) of \(H\) in \(\Gamma_0/V\)
is a finite \(\ell\)-group.
Indeed the maximal pro-\(\ell\) quotient of \(H\) is the quotient
by (TW3): normal cores show that using open normal subgroups
of \(\Gamma_0\), rather than all open normal subgroups of \(H\),
does not change it.
The images of \(W_N\) are contained in the iterated Frattini
subgroups of \(P\), which reach \(1\).
Equivalently, every nontrivial finite \(\ell\)-group has a
nonzero homomorphism to \(\mathbb F_\ell\), so each nontrivial
Frattini step strictly decreases its order.
This proves (TW4).

\paragraph{Choosing the primes, with the Frobenius convention fixed.}
Lift each \(s_j\) to \(g_j\in G_L\).
Since \(s_j\in G_0'\), its lift has
\(\varepsilon(g_j)=1\).
Let
\[
 U_N=\operatorname{pr}^{-1}(W_N)
                 \cap\ker(\varepsilon\bmod\ell^N)\subset G_L .
\]
This is open and normal.
Chebotarev, applied in the finite quotient by \(U_N\), chooses
distinct auxiliary places \(v_{Nj}\) of \(L^+\), outside \(S\),
split in \(L/L^+\), and places \(w_{Nj}\) of \(L\) above them,
with arithmetic Frobenius representatives \(\sigma_{Nj}\)
satisfying
\[
                         \sigma_{Nj}g_j\in U_N .
 \tag{TW5}
\]
Here one may choose degree-one primes of \(L\).
The usual avoidance of the finitely many rational primes
ramifying in \(L/L^+\), together with degree one, makes the
underlying \(L^+\)-place split; it also avoids every prescribed
finite set.
This is a finite Chebotarev problem at each \(N\), so simultaneous
choice of the \(r\) distinct places follows by finite induction.

Enlarge \(S\) by these \(v_{Nj}\) and let
\(h_N:\Gamma_N\twoheadrightarrow\Gamma_0\) be the resulting
arithmetic quotient map.
At \(w_{Nj}\), choose a topological generator \(\tau_{Nj}\)
of tame pro-\(\ell\) inertia and use the source convention
\[
             \phi_{Nj}=\operatorname{pr}_N(\sigma_{Nj}^{-1}),
 \qquad Q_{Nj}=\mathrm Nw_{Nj}.
 \tag{TW6}
\]
The inverse in this definition matters.
Cyclotomic compatibility with arithmetic Frobenius and (TW5)
give
\[
                           Q_{Nj}\equiv1\pmod{\ell^N}.
 \tag{TW7}
\]
Normality of \(W_N\), followed by inversion of (TW5), gives
\[
             h_N(\phi_{Nj})s_j^{-1}\in W_N .
 \tag{TW8}
\]
Consequently these products lie in \(H\) for every \(N\), and
converge to \(1\) modulo every open normal subgroup containing
\(\mathrm{PL}_\ell(H)\).

Put \(H_N=h_N^{-1}(H)\).
The tame relation, in the inverse-Frobenius convention (TW6),
is
\[
 \phi_{Nj}^{-1}\tau_{Nj}\phi_{Nj}
       =\tau_{Nj}^{Q_{Nj}}
                        \pmod{\mathrm{PL}_\ell(H_N)}.
 \tag{TW9}
\]
Wild inertia at this place has pro-residue-characteristic order,
hence has trivial image in a pro-\(\ell\) quotient.
The prime-to-\(\ell\) part of tame inertia also disappears.
The local kernel of the tame pro-\(\ell\) quotient maps into
\(\mathrm{PL}_\ell(H_N)\): every finite quotient in (TW3) has
an \(\ell\)-group image on \(H_N\), so it kills these local
prime-to-\(\ell\) subgroups.
The arithmetic kernel \( \ker h_N\) is the closed normal subgroup
generated by the new inertia groups.
It follows that
\[
 \begin{gathered}
 h_N(\tau_{Nj})=1,\\
 \ker h_N\subset
 \overline{\langle\!\langle
       \tau_{N1},\ldots,\tau_{Nr},
                         \mathrm{PL}_\ell(H_N)\rangle\!\rangle}.
 \tag{TW10}
 \end{gathered}
\]
The old local maps and their cyclotomic and polarization
identities are preserved by \(h_N\).

\paragraph{Uniform numbers of generators, without a Selmer assumption.}
There is a constant
\[
 b_H=
 \#\operatorname{Hom}_{\mathrm{cts}}(H,\mathbb F_\ell)\,
                   \ell^{\,r[\Gamma_0:H]}
 \tag{TW11}
\]
such that
\[
             \#\operatorname{Hom}_{\mathrm{cts}}
                         (H_N,\mathbb F_\ell)\leq b_H
                     \quad\hbox{for every }N.
 \tag{TW12}
\]
Here is the arithmetic kernel count.
Modulo homomorphisms inflated from \(H\), a character of \(H_N\)
is determined by its values at the \(r[\Gamma_0:H]\) elements
\(x\tau_{Nj}x^{-1}\), using one lift \(x\) of each coset
of \(H\) in \(\Gamma_0\).
Indeed conjugation by \(H_N\) does not affect a homomorphism
from \(H_N\) to an abelian group.
If all those values vanish, it kills the closed normal subgroup
generated by the new \(\tau_{Nj}\).
It also kills \(\mathrm{PL}_\ell(H_N)\), by (TW3).
Equation (TW10) then makes it descend to \(H\).
Each value has \(\ell\) choices, proving (TW11)--(TW12).
Thus the bound is obtained before either of the first- and
second-Selmer estimates.

\paragraph{Why the required points and first-order data descend.}
For a point of the fixed residual determinant, every
\(h\in H\) has residual characteristic polynomial \((X-1)^n\).
Its residual representation need not be the identity matrix.
Nevertheless its image on \(H\) modulo \(T^m\) is a finite
\(\ell\)-group: all residual elements are unipotent, so their
orders are powers of \(\ell\), and Cauchy's theorem shows that
their finite group has \(\ell\)-power order.
The congruence kernel is an \(\ell\)-group too.
It follows that the point kills \(\mathrm{PL}_\ell(H)\).
One may also see this directly from Cayley--Hamilton:
\[
 \bar\rho(h)\text{ unipotent}\quad\Longrightarrow\quad
     (\rho(h)^{\,\ell^{nm}}-1)\in T^mM_n(k[[T]]).
 \tag{TW13}
\]
Apply the same calculation to all conjugates of \(h\), take the
normal core of the finite representation kernel, and use (TW3).

This argument includes the odd generators and the actual
first- and second-order cocycles.
Their semidirect-product or square-zero-matrix representations
at finite coefficient precision have an additive
\(\ell\)-group kernel and the same residual unipotent image
on \(H\) or \(H_N\).
They therefore kill the corresponding \(\mathrm{PL}\).
The old local conditions are tested on the original local groups,
so composing their maps with the quotient does not change those
conditions.
Inflation consequently identifies the relevant finite Selmer
cocycles and boundaries on the original groups and on the
\(\mathrm{PL}\)-quotients.
In particular the first- and second-Selmer
bounds proved below transport with their original constants.

There is a corresponding statement for the completed determinant
and its paired Cayley--Hamilton algebra.
At any finite Artinian coefficient quotient the algebra is finite.
The residual characteristic polynomial of \(h\in H\) is
\((X-1)^n\), so \(h-1\) is nilpotent modulo the maximal ideal.
Nilpotence of that ideal makes \(h-1\) nilpotent in the finite
algebra.  A unit \(1+x\) with \(x\) nilpotent over a finite
\(\ell\)-power coefficient ring has \(\ell\)-power order,
by the binomial theorem.
Thus the finite \(H\)-image is again an \(\ell\)-group and
kills \(\mathrm{PL}_\ell(H)\).
Taking all finite coefficient quotients proves descent in the
completed topology, including the \emph{closed}
Cayley--Hamilton relations.
This is the same determinant and the same completed support
ideal used in the support-propagation subsection.
No new determinant ring is inserted at the limit step.

We may therefore pass to the actual quotient tower and henceforth
assume
\[
             \mathrm{PL}_\ell(H)=1,\qquad
             \mathrm{PL}_\ell(H_N)=1 .
 \tag{TW14}
\]
The sign and full cyclotomic character descend:
\(\eta|_H=1\), and the image of \(H\) in
\(1+\ell\mathbb Z_\ell\) is pro-\(\ell\).
Every map and every marked element above descends as well.
The image of \(H_N\) is now pro-\(\ell\), and (TW12) bounds its
minimal number of topological generators.
This is the Burnside basis theorem for pro-\(\ell\) groups:
generators of its finite Frattini quotient lift to topological
generators.  The coarse bound \(b_H\) is more than sufficient.

\paragraph{One finitely generated cover for every level.}
Let \(a_0=[\Gamma_0:H]\).
Choose \(a_0\) coset representatives of \(H\) and one additional
marked representative equal to complex conjugation \(c\).
Lift the coset representatives to each \(\Gamma_N\), using
\(c_N\) for the marked one, and add at most \(b_H\) generators
of \(H_N\).
The resulting \(a_0+1+b_H\) elements generate \(\Gamma_N\).
Let \(P\) be the free profinite group on this fixed number of
generators.
It has continuous surjections
\[
                       \pi_N:P\twoheadrightarrow\Gamma_N
 \tag{TW15}
\]
with one fixed \(c_P\in P\) mapping to \(c_N\), and with
\(h_N\pi_N\) independent of \(N\) after projection to
\(\Gamma_0/H\).
The kernel \(H_N\) is pro-\(\ell\) and finitely generated;
the finite coset representatives make its closure together with
these lifts the whole group.
This verifies surjectivity in (TW15), rather than merely
mapping a free group to each level.

\paragraph{The kernel limit and the base map.}
Choose an ultrafilter \(\mathcal U\) on \(\mathbb N\) containing
every cofinite set, and put \(K_N=\ker\pi_N\).
For every open normal \(V\subset P\), the images \(K_NV/V\)
range over the finite set of subgroups of \(P/V\).
Exactly one such image occurs on a set in \(\mathcal U\).
These images are compatible as \(V\) varies, and their transition
maps are surjective: this holds for the images of every \(K_N\)
and hence holds on the intersection of the finitely many
relevant sets in \(\mathcal U\).
Compactness produces a closed normal subgroup \(K_\infty\subset P\)
with exactly those finite images.
Set
\[
                        \Gamma_\infty=P/K_\infty .
\]
Then
\[
                K_NV=K_\infty V
                   \quad\text{for \(\mathcal U\)-almost all }N
 \tag{TW16}
\]
for every fixed \(V\).

Write \(f_N=h_N\pi_N\).
For any finite quotient \(\Gamma_0/V_0\), there are only finitely
many continuous homomorphisms from \(P\) to that quotient,
because \(P\) has finitely many topological generators.
Their \(\mathcal U\)-limits are compatible and define
a continuous \(f_\infty:P\to\Gamma_0\).
For each \(V_0\), the equality \(f_N=f_\infty\) modulo \(V_0\)
holds \emph{uniformly on all of \(P\)} for almost all \(N\).
This proves \(K_\infty\subset\ker f_\infty\), using (TW16).
It also proves surjectivity, since every \(f_N\) is surjective
onto every finite quotient and the image of a compact group
is closed.
We obtain
\[
                      h_\infty:\Gamma_\infty
                                      \twoheadrightarrow\Gamma_0 .
 \tag{TW17}
\]
The sign and cyclotomic characters and complex conjugation are
the ones obtained from \(f_\infty\).
In particular \(c_\infty^2=1\), since \(c_N^2=1\) at every level.

\paragraph{Uniform local limits; no lifting of local homomorphisms.}
Here an arbitrary pointwise limit is insufficient.
For each fixed old local group \(Y\), let
\(H_Y\) be the inverse image of \(H\) under its base map, and put
\(P_Y=\mathrm{PL}_\ell(H_Y)\), defined using \(Y\).
The local finiteness input and the pro-\(\ell\) Burnside basis
argument show that \(Y/P_Y\) is topologically finitely generated.
Every local map \(i_N:Y\to\Gamma_N\) kills \(P_Y\), by (TW14):
a finite quotient of \(H_N\) is an \(\ell\)-group, so the
pullback of each such quotient appears in the defining
intersection for \(P_Y\).

Fix \(V\triangleleft_{\mathrm{open}}P\).
For almost all \(N\), the subgroup \(S=K_NV=K_\infty V\)
is fixed.  At those levels there is a canonical quotient
\[
                         \Gamma_N\longrightarrow P/S ,
\]
because \(\ker\pi_N=K_N\subset S\).
Its composite with \(i_N\) is a continuous homomorphism
\(Y/P_Y\to P/S\).
There are finitely many such homomorphisms.
Thus it is a single \emph{whole local map}, simultaneously for
all \(y\in Y\), for almost all \(N\).
These stabilized maps are compatible as \(V\) varies, and give
a continuous local homomorphism
\(i_\infty:Y\to\Gamma_\infty\).
Its composite with \(h_\infty\) is the original base local map.

Equivalently, for every \(V\), the local elements at level \(N\)
and at the limit have lifts to \(P\) whose ratio lies in \(V\),
uniformly for all \(y\in Y\), for almost all \(N\).
This is the precise \srcid{Close} condition in the limit
definition.
The argument uses no continuous homomorphic lift
\(Y\to P\), which need not exist.
If one chooses arbitrary set sections of \(\pi_N\), their
reductions modulo \(S\) are the canonical quotient map just
described and are continuous; that is all that is required.
There are only finitely many old local places, so all these
conditions may be imposed simultaneously.

\paragraph{Marked Frobenius, inertia and the commuting limit.}
Choose set-theoretic lifts of \(\phi_{Nj}\) and \(\tau_{Nj}\)
to \(P\), and take their compact \(\mathcal U\)-limits.
Changing these lifts changes the limit only by \(K_\infty\),
by (TW16); hence they define elements
\(\phi_{\infty j},\tau_{\infty j}\in\Gamma_\infty\).
The finite set of marked elements has the same uniform
\srcid{Close} property.
Equations (TW4), (TW8), and (TW14) give
\[
 h_\infty(\phi_{\infty j})=s_j,\qquad
 h_\infty(\tau_{\infty j})=1 .
 \tag{TW18}
\]

Finally these two marked elements commute.
In the quotient tower, (TW9) is an exact identity.
At any fixed finite quotient arising from \(P/V\), the image of
\(\tau_{Nj}\) belongs to an \(\ell\)-group of bounded order,
because \(H_N\) is pro-\(\ell\).
Its exponent therefore divides some \(\ell^e\), independent of
\(N\) for that fixed finite quotient.
For \(N\geq e\), (TW7) makes
\(\tau_{Nj}^{Q_{Nj}}=\tau_{Nj}\) in that quotient.
The commutators consequently vanish in every finite quotient
at the limit:
\[
                        [\phi_{\infty j},\tau_{\infty j}]=1 .
 \tag{TW19}
\]
Together (TW15)--(TW19), the uniformly stabilized local maps,
and the cofinal ultrafilter are the actual limit data used in
the cohomology-limit and completed-support constructions.

\paragraph{Scope and exact frozen sources.}
This proves the arithmetic tower and the common group limit
from a fixed finite list (TW1).
The separable eigenvalues may be split by a finite extension of
the complete coefficient DVR.
For finitely many points, choose a common multiple of the residue
degrees, then extend their finite residue fields to the same
cardinality.  The resulting DVR extensions are finite, injective
and local; all point equations and odd generators extend with
them.
Thus the eigenvalue and common-residue-cardinality requirements
at the cohomology-limit interface are obtained by actual finite
coefficient extensions.

All paths below are under
\srcid{Development.AutomorphicArithmetic.AutomorphicVarieties}
at the frozen tag \srcid{alpoge-source-7cbf7a0c}.
The residual kernel is
\srcid{PatchingModules.ChebotarevFamilies.StarSingle.hEx};
its cofinal sequence is \srcid{FrSNormal.frat}.
The prime choice is \srcid{ChebOne} and
\srcid{ChebFam.chebFam}; the enlarged arithmetic quotient and
kernel are \srcid{LevUp.levUp} and \srcid{kerLe}.
The uniform character count is \srcid{HomBound.homBound},
and the residual-unipotent descent is \srcid{UnipPow.rhoPL}.
The assembled arithmetic construction is
\srcid{PatchingModules.AutomorphicDetection.ArPl.arPl}.
The quotient and cover are
\srcid{PatchingModules.SeparableMatrixStability.QuotEx.quot_ex},
\srcid{PlGen.pl_gen}, and \srcid{SysCover.sys_cover}.
The kernel limit uses \srcid{AlgebraTools.LimKer} and
\srcid{LimElt}; the uniform local argument is
\srcid{SeparableMatrixStability.Loc_fin.loc_fin} and
\srcid{InLim.in_lim}; the commuting limit is
\srcid{CommLim.comm_lim}.
Coefficient splitting and common residue degree are
\srcid{PatchingModules.DeterminantPadicFree.CharPOfPadicAlt.onePt}.
These sources are in
\url{https://github.com/ChristianKudeba/alpoge-rh-verification/tree/alpoge-source-7cbf7a0c/src/Development/AutomorphicArithmetic/AutomorphicVarieties}.



\subsection{The arithmetic second-Selmer bound}
\label{sec:arithmetic-second-selmer}

Here we prove the finite-level assertion \srcid{TWGrp.B2}, rather
than assuming it in the limit argument.
Let \(A=k[[T]]\), \(F=k((T))\), \(q=\#k\), with
\(\operatorname{char}k=\ell\ne2\).
Keep a nice paired point \(\rho\) of rank \(n\), with block sizes
\(a,b=n-a\), and the original odd Selmer generators
\(\kappa_1,\ldots,\kappa_g\).
The hypotheses on these objects used below are precisely:
\[
 \begin{gathered}
 \rho(G_0)\subset
   \mathrm{GL}_a(A)\times\mathrm{GL}_b(A),\qquad
 F[\rho(G_0')]=M_a(F)\oplus M_b(F),\\
 G_0'=\ker(\varepsilon:G_0\to\mathbb Z_\ell^\times),\\
 \psi_{v,i}\ne\psi_{v,j}\quad(i\ne j),\qquad
 T^s z-\sum_jc_j\kappa_j\text{ is an odd boundary}
 \text{ for every integral odd Selmer }z,
 \tag{B2a}
 \end{gathered}
\]
where \(s\) and \(c_j\in A\) in the last statement may depend on \(z\).
We also keep the actual polarized and unramified equations:
\(\rho(cgc^{-1})=\varepsilon(g)^{1-n}\rho(g)^{-\mathsf t}\),
\(\rho(I_u)=1\),
\(\kappa_j(cgc^{-1})=-\kappa_j(g)^{\mathsf t}\), and
\(\kappa_j(I_u)=0\).
Each \(\kappa_j\) is a continuous odd adjoint cocycle with its
integral ordinary local correcting matrices and prescribed
zero diagonal on the local \(\ell\)-primary torsion.
The ordinary bases are integral and respect the block type;
the characters \(\psi_{v,i}\) are trivial on that specified
local \(\ell\)-primary torsion.
The finite-coefficient Galois cohomology groups are finite.
We use the global and local Euler characteristics, local duality,
the global localization bound in degree two, and vanishing in degree
three for \(\ell\)-primary coefficients.  For the actual CM
arithmetic groups these are the cohomological statements proved
in the global-cohomology subsection: the group over the totally
real base has odd \(\ell\), and all primes above \(\ell\) are in the
ramification set.

Let \(G_N\to G_0\) be the actual auxiliary arithmetic groups.
For clarity, the Taylor--Wiles properties used in this proof are:
there is one fixed open normal \(H\subset\Gamma_0\) on which
\(\eta\), \(\bar\varepsilon\), and the residual character semisimplification
are trivial; there are fixed \(s_1,\ldots,s_r\in G_0'\), each
\(\rho(s_j)\) having distinct eigenvalues over \(\overline F\);
and there are auxiliary Frobenius and inertia elements
\(\phi_{Nj},\tau_{Nj}\) satisfying
\[
 \begin{gathered}
 h_N(\tau_{Nj})=1,\qquad
 \rho(h_N(\phi_{Nj}))\equiv\rho(s_j)\pmod{T^m}
       \quad(N\text{ sufficiently large for }m),\\
 \mathrm Nw_{Nj}\equiv1\pmod{\ell^N},\qquad
 \phi_{Nj}^{-1}\tau_{Nj}\phi_{Nj}
           =\tau_{Nj}^{\,\mathrm Nw_{Nj}}
               \pmod{\mathrm{PL}(H_N)},\\
 \ker h_N\subset
 \overline{\langle\!\langle\tau_{N1},\ldots,\tau_{Nr},
                      \mathrm{PL}(H_N)\rangle\!\rangle}.
 \tag{B2b}
 \end{gathered}
\]
Here \(\mathrm{PL}(H_N)\) is the kernel used to retain the maximal
pro-\(\ell\) quotient of the residual-kernel subgroup \(H_N\).
These are arithmetic prime-selection and tame inertia statements;
their construction is separate from the bound proved here.
The point kills \(\mathrm{PL}(H)\).
It suffices that every residual characteristic polynomial on \(H\)
is \((X-1)^n\); one need not assume the residual matrices are identity.
Their finite image consists of unipotent matrices and is therefore
an \(\ell\)-group: every element has \(\ell\)-power order and
Cauchy's theorem excludes other prime divisors of the group order.
The kernel of reduction modulo \(T\) is an \(\ell\)-group as well.
Thus every reduction modulo \(T^m\) has an \(\ell\)-group image on
\(H\), which kills \(\mathrm{PL}(H)\); separatedness proves the
integral assertion.

\paragraph{One uniform approximation lemma.}
We first record the complete-module argument responsible for uniformity
in \(m\).  Suppose \(C,C'\) are \(T\)-torsion-free \(A\)-modules,
\(C\) is complete and separated, \(C'\) is separated, and
\(\Phi:C\to C'\) is linear.  Put
\(S_m=\Phi^{-1}(T^mC')\).
Suppose \(S_1/TC\) is finite, and an \(A\)-submodule \(B\subset\ker\Phi\) has the
property that every \(z\in\ker\Phi\) has some \(T\)-power in \(B\).
Then, for one constant \(C_\Phi\) before all \(m\), the reduction
of \(S_m\) modulo \(T^mC\) is covered by at most \(q^{C_\Phi}\)
cosets of the reduction of \(B\).
The same is true after any linear projection of these reductions.

Here is a proof that does not assume \(C\) finitely generated.
The images of \(S_m\) in \(C/TC\) form a decreasing sequence of
subgroups of the finite image of \(S_1\); they stabilize, say from
\(m_0\) onwards.
Division by \(T\) is legitimate in \(C'\).
Here are the stabilization and limit steps explicitly.
For every \(j\geq1\), the images of \(S_m\) in \(C/T^jC\)
are constant for \(m\geq m_0+j-1\).
For the induction, let \(x\in S_{m_0+j}\) and choose a solution
\(s_1\) at any larger prescribed precision, with
\(x-s_1=T c_1\), using stabilization modulo \(T\).
Take that precision at least \(m_0+j\).
Then \(c_1\in S_{m_0+j-1}\), because
\(T\Phi(c_1)=\Phi(x-s_1)\) and \(C'\) is torsion free.
Apply the induction hypothesis to \(c_1\), approximating it
modulo \(T^j\) by a solution \(s_2\) at the prescribed precision.
The solution \(s_1+Ts_2\) approximates \(x\) modulo \(T^{j+1}\).
This proves the induction and, in particular, the convenient
slightly weaker assertion
\[
 \forall j,M,\quad
 x\in S_{m_0+j}\Longrightarrow
 \exists s\in S_M,\quad x-s\in T^jC.
 \tag{B2c1}
\]
Starting with \(r_0=x\in S_{m_0+j}\), apply (B2c1) successively
to choose
\[
 r_n\in S_{m_0+j+n},\qquad
 r_n-r_{n+1}=T^{j+n}d_n .
\]
Completeness gives \(L=\sum_{n\geq0}T^nd_n\in C\).
Set \(k=x-T^jL\).
For each \(n\), \(k-r_n\in T^{j+n}C\), whereas
\(\Phi(r_n)\in T^{m_0+j+n}C'\).
Separatedness of \(C'\) gives \(\Phi(k)=0\).
Moreover \(T^j\Phi(L)=\Phi(x)\in T^{m_0+j}C'\);
torsion freeness implies \(L\in S_{m_0}\).
Thus
\[
 S_{m_0+j}\subset\ker\Phi+T^jS_{m_0}\qquad(j\geq0).
 \tag{B2c}
\]
The kernel is complete and
\(\ker\Phi\cap TC=T\ker\Phi\), again by torsion freeness of \(C'\).
Thus its reduction modulo \(T\) is finite; successive approximation
over \(A\) proves it is a finite \(A\)-module.
The pointwise \(T\)-power condition on \(B\) is therefore uniform:
\(T^s\ker\Phi\subset B\) for some \(s\).
Choose finite representatives for \(\ker\Phi/B\) and for
\(S_{m_0}\) modulo \(T^{m_0}C\).
Formula (B2c) covers every \(S_{m_0+j}\) modulo \(T^{m_0+j}C\) by
the sums of one representative of the first kind and \(T^j\)
times one of the second.
There are uniformly finitely many such sums.
The finitely many \(m<m_0\) are absorbed into the constant.
Their finiteness follows inductively from finiteness at \(m=1\):
a difference divisible by \(T^r\) can be divided, and its new
residue is a solution modulo \(T\).
This proves the lemma.

Apply it to continuous integral global cochains together with their
finitely many local correcting matrices.
Let \(\Phi\) list the cocycle, polarization, inertia, ordinary
lower-entry and prescribed zero-diagonal defects.
Its target is a product of finite free \(A\)-modules.
The source is complete in the uniform \(T\)-adic topology:
compatible locally constant reductions define a continuous integral
cochain.  The solution set modulo \(T\) is finite by finite Galois
\(H^1\), finite boundaries, and finitely many local matrices.
Every finite-coefficient cochain lifts continuously, by using a
section of the finite coefficient reduction.
Take \(B\) to consist of exact solutions whose global cochain is
in the span of the \(\kappa_j\) and integral boundaries; leave the
compatible local correcting matrices unrestricted.
The last condition in (B2a) is exactly the \(T\)-power condition
on this \(B\).  Projecting to global cochains in the lemma yields
\[
 \#\operatorname{Sel}_{\mathrm{odd},0}(A/T^e)
                  \leq q^{ge+C_0}\qquad(e\geq0).
 \tag{B2d}
\]
This is a bound on classes, since the image of \(B\) has at most
\(q^{ge}\) classes.  It is not obtained merely by asserting that
the integral rank is \(g\); the approximation lemma controls
classes which exist only at finite precision.

\paragraph{The auxiliary primes add only a bounded odd Selmer factor.}
Let \(M_e=\operatorname{Odd}_a(A/T^e)\), with adjoint \(G_0\)-action,
extended to the full polarized group by minus transpose at complex
conjugation.
For each \(j\), the endomorphism
\[
 L_j:\operatorname{Odd}_a(A)\to\operatorname{Odd}_a(A),\qquad
       Y\longmapsto \rho(s_j)Y\rho(s_j)^{-1}-Y
\]
is injective after inverting \(T\).
Indeed the centralizer of a matrix with distinct eigenvalues is a
commutative algebra.  It contains the even block idempotent, so an
odd matrix in that centralizer commutes with the idempotent and is
zero.  Smith normal form consequently gives \(c_j\) with
\[
       \#\ker(L_j\bmod T^e)\leq q^{c_j}\quad\hbox{for every }e.
 \tag{B2e}
\]

For fixed \(m\) and sufficiently large \(N\), take an odd cocycle
modulo \(T^e\), \(e\leq m\).
Its value at \(\tau_{Nj}\) is unchanged by a coboundary, because
\(\rho(\tau_{Nj})=1\).
Its finite semidirect-product representation kills
\(\mathrm{PL}(H_N)\): the image of \(H_N\) in the residual kernel,
including the cocycle coordinate, is an \(\ell\)-group.
The tame relation (B2b) and
\(\mathrm Nw_{Nj}\equiv1\pmod\ell\) therefore show that its value
at \(\tau_{Nj}\) centralizes \(\rho(h_N(\phi_{Nj}))\).
This is the same centralizer as in (B2e) at precision \(e\).
If all these inertia values vanish, (B2b) makes the cocycle descend
to \(G_0\), with the same old local conditions.
Inflation on \(H^1\) is injective, so the inertia-value map gives
\[
 \#\operatorname{Sel}_{\mathrm{odd},N}(A/T^e)
 \leq \#\operatorname{Sel}_{\mathrm{odd},0}(A/T^e)
                                  \prod_jq^{c_j}
 \leq q^{ge+C_0+C_{\mathrm{aux}}},
 \quad C_{\mathrm{aux}}=\sum_jc_j.
 \tag{B2f}
\]
The constant \(C_{\mathrm{aux}}\) is fixed for the tower.

\paragraph{The actual relative obstruction group.}
At an ordinary place \(v\mid\ell\), let \(B_v\subset M_e\) be
the upper-triangular odd submodule in the chosen ordinary basis,
and put \(Q_v=M_e/B_v\).
At a prescribed non-\(\ell\) place put \(Q_u=M_e\), imposing
unramified classes on the specified inertia subgroup.
The ordinary type condition makes \(Q_v\) free of rank \(ab\)
over \(A/T^e\).
Define \(\mathcal H^2_e\) by pairs
\[
 \left(u\in Z^2(G_N,M_e),\
       b_w\in C^1(G_w,Q_w):\
       db_w=\pi_w\operatorname{res}_w u\right),
\]
modulo pairs coming from a global \(1\)-cochain and local
unramified \(1\)-cocycles.
At ordinary places unramified here means locally a boundary,
because the specified local subgroup is the full group.
This is precisely the group
\srcid{Dat.H2LocalConditions}; its defining cochains retain the
ordinary correcting matrices.

Write \(\mathcal C_e\) for the cokernel of
\(H^2(G_N,M_e)\to\prod_wH^2(G_w,Q_w)\).
The relative-cochain exact sequence, followed by Euler
characteristics, gives
\[
 \#\mathcal H^2_e
 \leq
 \#\operatorname{Sel}_{\mathrm{odd},N}(A/T^e)\,
 \left(\prod_{v\mid\ell}\#H^0(G_v,Q_v)\right)\#\mathcal C_e.
 \tag{B2g}
\]
Here is the cancellation explicitly.
The relative exact sequence has outer terms the cokernel of
\(H^1(G_N,M_e)\to\prod_wH^1(G_w,Q_w)/\mathcal L_w\)
and the kernel of localization in degree two.
Their cardinalities multiply to \(\#\mathcal H^2_e\).
The global odd module has rank \(2ab\), while minus transpose
has a fixed space of rank \(ab\) at each real place.
If \(d=[L^+:\mathbb Q]\), global Euler characteristic contributes
\(q^{-eabd}/\#H^0(G_N,M_e)\).
Local Euler characteristic at \(v\mid\ell\) contributes
\(q^{eab[L^+_v:\mathbb Q_\ell]}\#H^0(G_v,Q_v)\).
The degree sums are \(d\), so these powers cancel.
At the other local places, the unramified \(H^1\) cardinality
equals the \(H^0\) cardinality and contributes no remaining factor.
The finite prescribed inertia types have order prime to \(\ell\);
passing through them does not change that identity.
Discarding the positive global \(H^0\) denominator proves (B2g).

There are constants \(C_H,C_D\), independent of \(e,N\), with
\[
 \#H^0(G_v,Q_v)\leq q^{C_H},\qquad
 \#\mathcal C_e\leq q^{C_D}.
 \tag{B2h}
\]
For the first, choose for each distinct pair \(i,j\) an element
where \(\psi_{v,i}/\psi_{v,j}\ne1\).
In the lower-entry quotient, conjugation is triangular when entries
are ordered by their distance below the diagonal.
The diagonal coefficients of the invariance equations are these
nonzero character differences.
Successively eliminating the lower entries gives one nonzero
\(a_v\in A\) annihilating all invariants in every \(Q_v\).
Their cardinality is at most
\(q^{ab\,v_T(a_v)}\), uniformly in \(e\).
There are finitely many ordinary places.

For the second, global duality bounds \(\mathcal C_e\) by
\[
 \operatorname{Hom}_{\mathbb F_\ell[G_N]}(M_e,\mathbb F_\ell(1)).
\]
Localization first in the full module and then in the quotients
is legitimate because the local cohomological dimension is two.
Choose \(s\) with \(T^sE\in A[\rho(G_0')]\), where
\(E=\operatorname{diag}(1_a,0_b)\); this follows from the span
condition in (B2a).
An equivariant additive functional \(f\) satisfies
\(f(\rho(g)Y)=f(Y\rho(g))\) for \(g\in G_0'\), since the
cyclotomic character is \(1\) there.
The same holds with any \(A\)-multiple of \(Y\), so it extends to
the displayed \(A\)-linear span.
Thus
\[
                   f(T^s(EY-YE))=0 .
\]
On the upper and lower odd blocks the commutator is respectively
\(Y\) and \(-Y\).  Splitting the blocks proves
\(f(T^sM_e)=0\).
All such functionals are determined by \(M_e/T^sM_e\), and their
number is at most \(q^{2abs}\).
This proves the second assertion of (B2h), for example with
\(C_D=2abs\).

\paragraph{The obstruction quotient at every precision.}
A bound for \(\mathcal H^2_m\) alone is not enough.
For every \(e\leq m\), the two coefficient exact sequences
\[
 0\to A/T^{m-e}\xrightarrow{T^e}A/T^m\to A/T^e\to0,\qquad
 0\to A/T^e\xrightarrow{T^{m-e}}A/T^m\to A/T^{m-e}\to0
\]
give
\[
 \#(\mathcal H^2_m/T^e\mathcal H^2_m)
                         \leq\#\mathcal H^2_e\,\#\mathcal C_e.
 \tag{B2i}
\]
We spell out the maps and the defect argument.
Use the notation \(\mathcal E_r=\mathcal H^2_r\), and write
\[
 \begin{array}{c}
 f:\mathcal E_{m-e}\longrightarrow\mathcal E_m,\quad
 g:\mathcal E_m\longrightarrow\mathcal E_e,\\
 k:\mathcal E_e\longrightarrow\mathcal E_m,\quad
 r:\mathcal E_m\longrightarrow\mathcal E_{m-e}.
 \end{array}
\]
The first row is induced by the first coefficient sequence, the
second row by the second sequence.  The injections multiply
coefficients by \(T^e\) and \(T^{m-e}\), respectively, and
\[
                         f r=T^e .
 \tag{B2i1}
\]
The same coefficient maps act on the free lower-entry quotients
\(Q_w\), where both sequences are exact.
Local unramified cocycles lift under the coefficient surjections:
at an ordinary place they are boundaries, which lift by lifting
their correcting matrix; at another prescribed place first restrict to the finite
prime-to-\(\ell\) inertia quotient.
Averaging makes this restriction a boundary.
Lift its correcting matrix, subtract that boundary, and lift
the remaining Frobenius value in the inertia-invariant module;
averaging makes invariants exact under the coefficient
surjection.
A cocycle on the procyclic unramified quotient is specified
by that Frobenius value.
It is continuous for finite \(\ell\)-primary coefficients:
after an exponent of the finite Frobenius action, a further
\(\ell\)-power makes the norm sum zero.
Add back the lifted boundary.
This gives an actual lift of the original unramified cocycle,
including its finite inertial part.

It follows directly on representatives that
\(\ker g\subset\operatorname{im}f\).
Indeed, if a relative pair \(p\) maps to a relative coboundary,
lift its global correcting \(1\)-cochain and its local unramified
\(1\)-cocycles.
Subtract that lifted relative coboundary from \(p\).
Both its global and local coordinates now map to zero pointwise.
Exactness of the coefficient sequences gives unique preimages
under the first injection.  Injectivity makes those preimages
a global \(2\)-cocycle and compatible local \(1\)-cochains.
They define a class of \(\mathcal E_{m-e}\) mapping to \([p]\).
Consequently
\[
                 \#(\mathcal E_m/\operatorname{im}f)
                                    \leq\#\mathcal E_e.
 \tag{B2i2}
\]

For the other factor we give the actual defect construction.
Take a relative representative
\(p'=(u',b'_w)\) over \(A/T^{m-e}\).
Global degree-three vanishing makes
\(H^2(M_m)\to H^2(M_{m-e})\) surjective.
Hence \(u'\) lifts to an actual global \(2\)-cocycle \(u\):
first lift its cohomology class, then lift the correcting
\(1\)-cochain and subtract its differential.
Lift the local \(1\)-cochains \(b'_w\) arbitrarily to \(b_w\).
The local defects have unique expressions
\[
 \pi_w\operatorname{res}_w u-db_w
               =T^{m-e}c_w,\qquad c_w\in Z^2(G_w,Q_{w,e}).
 \tag{B2i3}
\]
The right side is interpreted using the coefficient injection
\(Q_{w,e}\hookrightarrow Q_{w,m}\).
Its injectivity and \(d^2=0\) prove that \(c_w\) is a cocycle.
The defect is measured in the \emph{actual} cokernel
\[
 \mathcal C_e=
  \left(\prod_wH^2(G_w,Q_{w,e})\right)/
       \operatorname{im}\!\left(H^2(G_N,M_e)
                       \longrightarrow\prod_wH^2(G_w,Q_{w,e})\right).
\]

The zero criterion is explicit.
If \(([c_w])_w\) is the localization of a global class represented
by \(v\in Z^2(G_N,M_e)\), choose local \(a_w\) such that
\(c_w=\pi_w\operatorname{res}_wv-da_w\).
Then
\[
       (u-T^{m-e}v,\ b_w-T^{m-e}a_w)
\]
is a relative cocycle over \(A/T^m\) mapping to \(p'\).
The same calculation applies to a difference of two relative
representatives.  Choose one representative and one system of
lifts for each element of
\(\mathcal E_{m-e}/\operatorname{im}r\).
Sending it to its defect class in \(\mathcal C_e\) is an
injective map of sets: equality of two defect classes makes
their difference a lift by the displayed correction, so the
two quotient classes agree.
This formulation needs no unproved independence of arbitrary
choices of relative representatives.  In particular,
\[
           \#(\mathcal E_{m-e}/\operatorname{im}r)
                                            \leq\#\mathcal C_e.
 \tag{B2i4}
\]
By (B2i1), \(f\) induces a surjection from this quotient onto
\(\operatorname{im}f/T^e\mathcal E_m\).
Multiplying the bounds for
\(\mathcal E_m/\operatorname{im}f\) and
\(\operatorname{im}f/T^e\mathcal E_m\), using (B2i2)--(B2i4),
proves (B2i).

Combining (B2f)--(B2i), with \(r_\ell\) the number of ordinary
places, yields, simultaneously for every \(e\leq m\) and all
sufficiently large \(N\),
\[
 \#(\mathcal H^2_m/T^e\mathcal H^2_m)
           \leq q^{ge+C},\qquad
 C=C_0+C_{\mathrm{aux}}+r_\ell C_H+2C_D.
 \tag{B2j}
\]
The constant is fixed before \(m,N\).

\paragraph{Primitive relations and the explicit second-order lift.}
Write the finite \(A\)-module \(\mathcal H^2_m\) in elementary
divisors, with lengths
\(\lambda_1\geq\lambda_2\geq\cdots\).
Equation (B2j) says
\(\sum_i\min(e,\lambda_i)\leq ge+C\).
Taking \(e=\lambda_{g+1}\), when it is nonzero, proves
\(\lambda_{g+1}\leq C\).
Thus \(T^C\mathcal H^2_m\) is generated by at most \(g\) elements.
For any \(g+1\) elements \(\xi_i\) of \(\mathcal H^2_m\), lift
their coordinates in such a generating set to a \(g\) by \(g+1\)
matrix over \(A\).  Its kernel over \(F\) has a nonzero vector.
Clear denominators and divide by its common \(T\)-power.
The resulting \(c_i\in A\) include a unit and satisfy
\[
                         T^C\sum_i c_i\xi_i=0 .
 \tag{B2k}
\]

For the even Selmer tuples \(z_{ij}\) in \srcid{B2}, use the
actual square-zero coefficient algebra
\[
 \mathbb D_g(A/T^m)
  =(A/T^m)\oplus\bigoplus_j\epsilon_j(A/T^m)
    \oplus\bigoplus_j\epsilon'_j(A/T^m)\oplus\eta(A/T^m),
 \quad\epsilon_i\epsilon'_j=\delta_{ij}\eta.
\]
The fixed first-order representation is
\((1+\sum_j\epsilon_j\kappa_j)\rho\).
The extension of coefficient modules governing the next step is
\[
 0\to\eta\,\operatorname{Odd}_a(A/T^m)
       \longrightarrow W_2
       \longrightarrow
       \bigoplus_j\epsilon'_j\operatorname{Even}_a(A/T^m)\to0.
\]
For the \(i\)-th tuple its connecting obstruction is the relative
class \(\xi_i\) just defined.
Its global two-cocycle is, up to the sign convention for \(d^1\),
\[
 (g,h)\longmapsto
        \sum_j[\kappa_j(g),\operatorname{Ad}(\rho(g))z_{ij}(h)].
 \tag{B2l}
\]
This formula follows immediately by multiplying the two
\(\mathbb D_g\)-valued matrices.
The local component is obtained by applying the same differential
to the chosen even local correcting matrices, using the fixed
odd correcting matrices for the \(\kappa_j\).
Its lower-entry projection has differential equal to the
restriction of (B2l), so it is a pair defining \(\mathcal H^2_m\).
This is a cochain construction, not an assumed obstruction
vanishing theorem.

When a linear combination of the relative classes vanishes,
a global odd one-cochain and local odd correcting matrices
trivialize the combination.
Putting that cochain in the \(\eta\)-coordinate gives
\[
 \Psi(g)=
    \sum_j\epsilon'_j\!\left(T^C\sum_i c_i z_{ij}(g)\right)
                           +\eta\,y(g).
\]
The cocycle equation is exactly cancellation of (B2l).
Complex-conjugation compatibility is obtained by averaging with
its polarized transform; \(2\) is a unit.
At ordinary places the same local correcting matrices make every
lower entry zero.  The required diagonal entries on the specified
inertia also remain zero: in the ordinary basis the relevant
commutators are commutators of upper-triangular matrices and have
zero diagonal.
At other prescribed inertia the relative unramified condition
makes the cocycle a boundary there.  Its action is trivial, by
\(\rho(I_u)=1\) and \(\kappa_j(I_u)=0\), so that boundary is zero.
Thus \(\Psi\in\operatorname{Sel}_{2,N}(A/T^m)\) and
\[
       \epsilon'_j(\Psi)=T^C\sum_i c_i z_{ij},\qquad
                    \text{some }c_i\in A^\times .
\]
This proves \srcid{TWGrp.B2} with one \(C\) before every
coefficient precision and every tower level.

\paragraph{Scope and exact sources.}
The proof applies in every rank and to both paired blocks; for
\(a=n\) the odd module is zero and the same argument simplifies.
It proves the actual finite-level arithmetic second-Selmer bound
from the displayed nice-point, original-generator, standard
arithmetic cohomology, and Taylor--Wiles group properties.
It does not silently manufacture the regular cosets or the
Taylor--Wiles tower; these are the separate prime-selection
inputs in (B2b).

All source paths below are at the frozen tag
\srcid{alpoge-source-7cbf7a0c}, under
\srcid{Development.AutomorphicArithmetic.AutomorphicVarieties}.
The uniform approximation proof is
\srcid{PatchingModules.DegreeCountCohomology.SolCnt}
(\srcid{scS\_stab}, \srcid{scS\_lim}, \srcid{sc\_kerfin},
\srcid{sc\_unif}, \srcid{solCnt}).
Its projected original odd-Selmer application is
\srcid{ArchimedeanH0Cards.ImgCard.sel\_cnt} and \srcid{sel0}.
The inertia comparison is
\srcid{ArchimedeanH0Cards.ExistsCore.odd\_kill},
\srcid{fix\_count}, and \srcid{selN}.
The local invariant bound and dual invariant bound are
\srcid{H0LocalCardBoundedOfClsKerLeKer} and \srcid{DualBound}.
The relative group is defined in
\srcid{PushforwardPreservesSelmerCocycles.Dat.H2LocalConditions};
its cardinal formula is \srcid{OdCount.odCount}.
The two exact-sequence chase is
\srcid{ObstructionBoundOfDuality.step}, using
\srcid{ExactMid} and \srcid{CokerBound}; primitive relations are
\srcid{GenTrick.genTrick}.
The explicit lift from those relations is
\srcid{Transl.transl}; the final arithmetic specialization is
\srcid{AutomorphicDetection.TaylorWilesObstructionBound.b2Ar}.
These file paths are below
\url{https://github.com/ChristianKudeba/alpoge-rh-verification/tree/alpoge-source-7cbf7a0c/src/Development/AutomorphicArithmetic/AutomorphicVarieties}.

\subsection{Trace detection and the arithmetic first-Selmer bound}
\label{sec:arithmetic-first-selmer}

We prove the actual first-order bound \srcid{TWGrp.B1}.
In particular the finite list of Taylor--Wiles elements is
constructed, rather than postulated to detect all dual classes.

Write \(A=k[[T]]\), \(F=k((T))\), \(Q=\#k\), and
\(\operatorname{char}k=\ell\ne2\).
Let \(\rho:G_0\to\mathrm{GL}_n(A)\) be one of the nice paired
points, with block sizes \(a,b=n-a\).
Put
\[
 E=M_a(F)\oplus M_b(F),\qquad
 O=\operatorname{Hom}(F^b,F^a)\oplus
                         \operatorname{Hom}(F^a,F^b).
\]
The precise global hypotheses used here are
\[
 F[\rho(G_0')]=E,\qquad
 G_0'=\ker(\varepsilon|_{G_0}),\qquad
 \exists g_*\in G_0:\bar\varepsilon(g_*)\ne1,
 \tag{B1a}
\]
together with polarization and the original odd generators
\(\kappa_1,\ldots,\kappa_g\).
Their independence means that a linear combination which is
the boundary of an odd skew matrix has all coefficients zero.
Their original spanning property is the one already used in
the second-Selmer proof.
At \(v\mid\ell\), the integral ordinary bases respect the block
type, and the ordered characters satisfy
\[
 \psi_{v,i}\ne\psi_{v,j}\ (i\ne j),\qquad
 \psi_{v,i}\ne\psi_{v,j}\bar\varepsilon_v\ (i<j).
 \tag{B1b}
\]
The prescribed local torsion is killed by the diagonal
characters.  At the other prescribed inertia groups,
\(\rho=1\) and \(\kappa_j=0\).
We use the finite arithmetic cohomology, local and global duality,
and Euler characteristics established previously.

Let \(D_0\subset G_0\) be the distinguished local image.
The remaining nice-point hypothesis used in this bound is exactly
\[
 \{Y\in M_n(F):
    \bar\varepsilon(g)\rho(g)Y\rho(g)^{-1}=Y
                         \text{ for all }g\in D_0\}=0 .
 \tag{B1c}
\]
This is \srcid{N4}, an explicitly specified local condition.
We do not replace it by irreducibility of \(\rho\).

\paragraph{The actual dual coefficient module and its traces.}
Use the square-zero coefficient algebra
\[
 {\mathbb D}_g=A\oplus\bigoplus_j\epsilon_jA
                   \oplus\bigoplus_j\epsilon'_jA\oplus\eta A,
 \qquad \epsilon_i\epsilon'_j=\delta_{ij}\eta,
\]
with all remaining products of positive-degree generators zero.
The fixed first-order representation is
\[
                   \rho_1=(1+\sum_j\epsilon_j\kappa_j)\rho .
\]
The primal and dual modules are
\[
 W_1=E_A\oplus\bigoplus_j\epsilon_jO_A,\qquad
 W_1^\vee=\eta E_A\oplus\bigoplus_j\epsilon'_jO_A.
 \tag{B1d}
\]
Their pairing is the \(\eta\)-coefficient of \(\operatorname{tr}(XY)\).
It is perfect over \(A\).
The \(G_0\)-action on the dual module is
\(\bar\varepsilon\,\operatorname{Ad}\rho_1\); its extension across
complex conjugation is minus the cyclotomic multiplier times
transpose.
For a continuous polarized twisted cocycle \(z\), put
\[
             \ell_s(z)=
       \eta\operatorname{tr}\bigl(z(s)\rho_1(s)\bigr)
                         \quad(s\in G_0').
 \tag{B1e}
\]
These are the actual trace functionals used by the construction.

We first prove that their common zero space over \(F\) consists
of boundaries.
Write the coordinates of \(z\) as
\[
                 z(g)=\eta\varphi(g)
                              +\sum_j\epsilon'_j\lambda_j(g),
 \qquad \chi=\bar\varepsilon .
\]
Then
\[
 \begin{split}
 \lambda_j(gh)&=\lambda_j(g)+
                  \chi(g)\operatorname{Ad}\rho(g)\lambda_j(h),\\
 \varphi(gh)&=\varphi(g)+\chi(g)\left(
       \operatorname{Ad}\rho(g)\varphi(h)+
       \sum_j[\kappa_j(g),
                     \operatorname{Ad}\rho(g)\lambda_j(h)]\right).
 \end{split}
 \tag{B1f}
\]
The odd \(\lambda_j\) and even \(\varphi\) have the same transpose
polarization sign \(s\), with \(s^2=1\).
For \(g\in G_0'\), vanishing of (B1e) is exactly
\[
       \operatorname{tr}\left(
        \varphi(g)\rho(g)+
                \sum_j\lambda_j(g)\kappa_j(g)\rho(g)\right)=0 .
 \tag{B1g}
\]

\paragraph{A relation calculation that detects the odd coordinates.}
Here are the algebraic steps behind this assertion.
Let \(R:F[G_0']\to E\) be induced by \(\rho\), and for
\(c=\sum c_g[g]\) write
\[
 K_j(c)=\sum_gc_g\kappa_j(g)\rho(g),\qquad
 L_j(c)=\sum_gc_g\lambda_j(g)\rho(g).
\]
The cocycle equations make \(K_j\) and \(L_j\) derivations
relative to \(R\) on this group algebra.
Define also
\[
 P(c)=\sum_gc_g\left(\varphi(g)\rho(g)+
                     \sum_j\lambda_j(g)\kappa_j(g)\rho(g)\right).
\]
Multiplying (B1f) gives the explicit product formula
\[
 P(cd)=P(c)R(d)+R(c)P(d)+
                  \sum_j\bigl(K_j(c)L_j(d)+L_j(c)K_j(d)\bigr).
 \tag{B1h}
\]
We have \(\operatorname{tr}P(c)=0\) by (B1g).
For \(c,d\in\ker R\) and opposite-block indices \(i,j\), insert
preimages under \(R\) of the two diagonal matrix units
\(e_{jj}\) and \(e_{ii}\) in \(\operatorname{tr}P(c e_{jj}d e_{ii})\).
Surjectivity of \(R\), which is (B1a), permits these choices.
Every term involving \(R(c)\) or \(R(d)\) vanishes.
Formula (B1h) leaves precisely
\[
 \sum_\nu\left(
      K_\nu(c)_{ij}L_\nu(d)_{ji}
                 +L_\nu(c)_{ij}K_\nu(d)_{ji}\right)=0 .
 \tag{B1i}
\]

We need the following elementary consequence of relation vanishing.
Suppose \(K\) is an untwisted odd cocycle and \(L\) is a
\(\chi\)-twisted odd cocycle.  If
\[
       \left(\sum_gc_g(K(g)+L(g))\rho(g)\right)_{ij}=0
                                  \quad(c\in\ker R)
 \tag{B1j}
\]
for one opposite-block pair \(i,j\), their corresponding entire
corner is principal on \(G_0\).
First, multiply the relation on the left and right by images
of elements of \(G_0'\).
The cocycle equation shows that the extra terms contain
\(R(c)=0\); hence (B1j) is preserved.
The matrix units in \(E\) move the indices to any entry in that
corner.
Thus the corner of \((K+L)(g)\rho(g)\) factors through a derivation
of \(E\) into that corner bimodule.

Every derivation of \(E=M_a(F)\oplus M_b(F)\) into \(M_n(F)\)
is inner.  An explicit formula is useful here.
Choose in each nonempty block an anchor index \(\beta(i)\),
put \(u_i=e_{i,\beta(i)}\), \(v_i=e_{\beta(i),i}\), and for a
derivation \(d\) take \(Y=\sum_i d(u_i)v_i\).
The identities
\(\sum_i u_iv_i=1\) and the Leibniz rule give
\(d(e_{ij})=Ye_{ij}-e_{ij}Y\) whenever \(i,j\) are in one block.
Linearity proves the assertion for all of \(E\).
After changing the sign of \(Y\), we obtain
\[
                  K+L=\delta Y\quad\hbox{on }G_0'
\]
for the chosen corner.

The two summands separate because \(\chi(g_*)\ne1\).
For completeness, if \(c_*=\chi(g_*)\), direct expansion of
the cocycle equations at \(g g_*\) and \(g_*g_*^{-1}gg_*\)
gives
\[
 \begin{split}
 Y_1=(c_*-1)^{-1}\bigl(
  c_*\operatorname{Ad}\rho(g_*)Y
             -c_*K(g_*)-L(g_*)-Y\bigr),\\
             K=\delta Y_1,\qquad L=\delta(Y-Y_1)
                                \quad\hbox{on }G_0'.
 \end{split}
 \tag{B1k}
\]
Here \(K,L\) denote their chosen corner projections.
Project the correcting matrices to that corner.
The errors from these equations, extended to \(G_0\) with their
respective twists, vanish on \(G_0'\).
Their values therefore centralize \(\rho(G_0')\).
An odd matrix centralizing \(E\) is zero, because it commutes
with the two block idempotents.
Thus the corner equations hold on all of \(G_0\).

Polarization reconstructs the opposite corner.
If \(K(\sigma g)=-K(g)^{\mathsf t}\), replace its corner
correcting matrix \(Y\) by \(Y-Y^{\mathsf t}\).
If \(L(\sigma g)=sL(g)^{\mathsf t}\), replace its corner
correcting matrix by \(Y+sY^{\mathsf t}\).
The identity
\[
        \delta_\chi(Y)(\sigma g)
                        =\delta_\chi(Y^{\mathsf t})(g)^{\mathsf t}
\]
follows from the polarized equation for \(\rho\).
These replacements give the full odd boundaries with the
required signs.

Apply this consequence to
\(K=\sum_\nu\alpha_\nu\kappa_\nu\),
\(L=\sum_\nu\beta_\nu\lambda_\nu\).
If (B1j) holds, the preceding argument makes \(K\) an odd skew
boundary; independence of the \(\kappa_\nu\) implies
\(\alpha_\nu=0\) for every \(\nu\).
In (B1i), fix \(d\in\ker R\), and put
\[
                 \alpha_\nu=L_\nu(d)_{ji},\qquad
                 \beta_\nu=K_\nu(d)_{ji}.
\]
Equation (B1i) is exactly (B1j), so all
\(L_\nu(d)_{ji}\) vanish.
Now each \(\lambda_\nu\) kills relations in the opposite entry.
Applying the same argument with \(\alpha=0\) and
\(\beta_\nu=1\), the other \(\beta\)'s zero, proves
\[
        \lambda_\nu=\delta_\chi Y_\nu,\qquad
                Y_\nu\in O,\quad sY_\nu^{\mathsf t}=Y_\nu .
 \tag{B1l}
\]
If \(a=n\), the odd module is zero and this paragraph is vacuous.

\paragraph{The even coordinate and the full dual boundary.}
Subtract from \(\varphi(g)\) the explicit cross term
\[
             \chi(g)\sum_\nu
                  [\kappa_\nu(g),\operatorname{Ad}\rho(g)Y_\nu].
\]
Call the result \(\varphi_0(g)\).
Equations (B1f) and (B1l) show, by expansion, that it is an
ordinary \(\chi\)-twisted even cocycle.
Cyclicity of trace and (B1g) give
\(\operatorname{tr}(\varphi_0(g)\rho(g))=0\) on \(G_0'\).
For a relation \(c\in\ker R\), put
\(M=\sum c_g\varphi_0(g)\rho(g)\in E\).
Apply the trace equation to \(gh\) and sum over \(c\).
It gives \(\operatorname{tr}(M\rho(h))=0\) for all \(h\in G_0'\).
The trace pairing on \(E\) is nondegenerate, so \(M=0\).
Thus \(\varphi_0(g)\rho(g)\) factors through an even derivation
of \(E\), which is inner by the same matrix-unit calculation.

Subtract the resulting even boundary.
The error vanishes on \(G_0'\), so its values centralize \(E\)
and therefore commute with every \(\rho(g)\).
The quotient \(G_0/G_0'\) is abelian.
Its two cocycle expansions at \(g g_*\) and \(g_*g\) yield
\[
      (c_*-1)\varphi_{\rm err}(g)
                =(\chi(g)-1)\varphi_{\rm err}(g_*).
\]
This is the boundary of
\((c_*-1)^{-1}\varphi_{\rm err}(g_*)\).
Averaging the correcting matrix with \(s\) times its transpose
makes it have the required polarization, since \(2\) is a unit.
Together with (B1l), this produces
\[
 Y=\eta Y_0+\sum_\nu\epsilon'_\nu Y_\nu\in W_1^\vee\otimes_A F,
 \qquad
             z(g)=\chi(g)\operatorname{Ad}\rho_1(g)Y-Y .
 \tag{B1m}
\]
Clearing the finitely many matrix denominators proves that an
integral \(z\) annihilated by all \(\ell_s\), \(s\in G_0'\),
has a \(T\)-power in the integral polarized boundary module.
This is the precise assertion called \srcid{Detects}.

\paragraph{Choosing finitely many simultaneous regular detectors.}
The module of continuous twisted cocycles with values in a fixed
finite free \(A\)-module is finite over \(A\).
Indeed its reduction modulo \(T\) injects into the finite set
of residual continuous cocycles: the kernel of reduction is
\(T\) times the integral cocycle module.
Dividing a cocycle by \(T\) preserves continuity by testing
one higher coefficient precision.
The module is complete in the uniform \(T\)-adic topology,
so successive approximation, or topological Nakayama, proves
finite generation.
Residual cocycle finiteness follows from \srcid{PhiCond}:
on the open kernel of the residual action a cocycle is a
homomorphism into a finite elementary \(\ell\)-group, and
there are finitely many choices on its finite complement.
This argument applies to the full matrix coefficient module
over \(\mathbb D_g\), and hence to all its submodules.

For finitely many nice points, let \(\mathcal S\subset G_0'\)
be the set where every point has separable characteristic
polynomial.
For each point this regular set is open: the discriminant is
a continuous \(A\)-valued function, and its nonvanishing is
open by separatedness of \(A\).
The actual regular-coset property \srcid{N3} says this open
set meets every nonempty relatively open part of \(G_0'\).
Finite induction, retaining the regular open sets already
obtained, shows that \(\mathcal S\) is dense in \(G_0'\).

For each integral cocycle \(z\), the function
\(s\mapsto\ell_s(z)\) is continuous.
Hence vanishing on \(\mathcal S\) implies vanishing on all
of \(G_0'\), and (B1m) applies.
On the finite \(A\)-module of integral cocycles, the functionals
\(\ell_s\), \(s\in\mathcal S\), span a finite submodule of its
\(A\)-dual.
Select finitely many which generate that span.
Their common kernel is therefore the common kernel of all
the functionals.
Do this for every point and take the finite union.
It produces one list \(s_1,\ldots,s_r\in\mathcal S\) which
detects every point, all of whose characteristic polynomials
are separable at every point.
The finite DVR extensions described in the tower subsection
split all these polynomials and give the same residue-field
cardinality across the finite family.
This proves the actual choice input.

\paragraph{A uniform bound on the dual Selmer factor.}
The trace calculation over \(F\) alone does not bound classes
which exist only modulo \(T^m\).
Apply the uniform approximation lemma from the second-Selmer
proof to continuous cochains, with defect map listing
their cocycle, polarization, and the finitely many
\(\ell_{s_j}\) equations.
Take \(B\) to be the \(A\)-submodule of polarized boundaries.
The hypothesis that every exact solution has a \(T\)-power
in \(B\) is (B1m); residual solutions are finite by
\srcid{PhiCond}.
The lemma consequently gives one \(C_{\rm tr}\), before all \(m\),
such that
\[
 \#\{z\in\operatorname{TwZ}(A/T^m):
                     \ell_{s_j}(z)=0\ \forall j\}
       \leq \#\operatorname{TwB}(A/T^m)\,Q^{C_{\rm tr}} .
 \tag{B1n}
\]

Use the actual arithmetic tower constructed in this section.
For fixed \(m\) and all sufficiently large \(N\), a full
polarized dual cocycle locally a boundary at each new place
has value zero at \(\tau_{Nj}\), because \(\tau_{Nj}\) acts
trivially on the coefficient module.
Its finite semidirect-product image kills
\(\mathrm{PL}_\ell(H_N)\), as proved in the tower construction.
The normal-kernel assertion (TW10) makes it descend to
\(\Gamma_0\).
Normalize its value at complex conjugation to zero by a
boundary, using \(2^{-1}\), and restrict to \(G_0\).

At finite precision there are only finitely many such
base cocycles, not just finitely many point matrices.
Intersect their finite semidirect-product kernels, along
with the kernels of \(\rho_1\) and the cyclotomic character
at that precision.
This open subgroup contains \(\mathrm{PL}_\ell(H)\).
Frobenius convergence (TW8) therefore identifies their
values at \(h_N(\phi_{Nj})\) and at \(s_j\), simultaneously.
The boundary equation at \(\phi_{Nj}\), multiplied by
\(\rho_1(s_j)\) and traced, gives
\[
 \eta\operatorname{tr}\bigl(
   (\operatorname{Ad}\rho_1(s_j)Y-Y)\rho_1(s_j)\bigr)=0 .
\]
Thus (B1n) applies.
Normalization at complex conjugation has exactly the same
fibers for cocycles and boundaries; canceling those fibers
gives the actual dual bound
\[
       \#H^1_{\mathcal L^\perp,N}
                    (W_1^\vee/T^m)\leq Q^{C_{\rm tr}},
 \tag{B1o}
\]
when at the auxiliary places the dual local condition is
being a boundary.
For the global counting inequality it suffices to impose
only these auxiliary local conditions, so any further
orthogonality conditions can only decrease this cardinality.

\paragraph{The local ordinary and invariant counts.}
Put
\[
 \nu=\frac{a(a-1)}2+\frac{b(b-1)}2+gab .
\]
Then \(W_1\) has rank \(n+2\nu\), and its upper-triangular
submodule in the actual ordinary basis has rank \(n+\nu\).
For an ordinary place \(v\), let
\(d_v=[L^+_v:\mathbb Q_\ell]\), and let \(d'_v\) be the free
\(\mathbb Z_\ell\)-rank of the local unit quotient used by the
ordinary weight algebra.
Local class field theory gives
\[
                         \sum_{v\mid\ell}d_v
                        =\sum_{v\mid\ell}d'_v=d,
             \qquad d=[L^+:\mathbb Q].
 \tag{B1p}
\]
The local abelian quotient, with the prescribed finite torsion
removed, has rank \(d'_v+1\), including its unramified factor.

Let \(Z_{b,v,m}\) be the actual upper-triangular cocycles,
with diagonal zero on that prescribed torsion.
Their diagonal coordinates are \(n\) continuous characters
of this abelian quotient.
They have at most \(Q^{mn(d'_v+1)}\) choices.
Filter the strictly upper-triangular coordinates by their
distance from the diagonal, retaining the even block entries
and the \(g\) odd copies.
There are \(\nu\) root coordinates.
Each associated rank-one character is
\(\psi_{v,i}/\psi_{v,j}\), \(i<j\).
Its degree-two local cohomology is uniformly bounded in \(m\):
local duality and (B1b) identify it with the dual of invariants
of a nontrivial character.
A fixed nonzero character difference annihilates those
invariants at every precision.
Local Euler characteristic, together with the elementary
boundary count, therefore bounds its cocycles by
\(Q^{m(d_v+1)+C_{ij}}\).
Successively choosing the filtered coordinates never gives
more choices than the product of these cocycle counts.
Thus, with \(C_v\) fixed before \(m\),
\[
 \#Z_{b,v,m}\leq
       Q^{\,m[n(d'_v+1)+\nu(d_v+1)]+C_v}.
 \tag{B1q}
\]
This counts the actual ordinary cocycles; the fixed local
correctors in \(\rho_1\) enter by the integral basis change.

There are also constants \(C_G,C_*,C_j\) such that
\[
 \begin{array}{rl}
 \#H^0(\Gamma_N,W_1/T^m)&\leq Q^{C_G},\\
 \#H^0(D_0,W_1^\vee/T^m)&\leq Q^{C_*},\\
 \#H^0(\langle s_j\rangle,W_1^\vee/T^m)
                                      &\leq Q^{mn+C_j}.
 \end{array}
 \tag{B1r}
\]
These assertions follow by finite matrix equations and Smith
normal form, as follows.
A \(G_0'\)-invariant constant even matrix centralizes \(E\),
hence is scalar on each block.
The minus-transpose equation at complex conjugation kills
both scalars.
The remaining odd components centralize \(E\) and vanish.
Thus the first invariant space is zero after inverting \(T\).
For the second, first apply (B1c) to every \(\epsilon'\)-coordinate,
then to the remaining \(\eta\)-coordinate; the cross terms
disappear, so its invariant space also vanishes over \(F\).
For the third, the odd components centralizing the regular
\(\rho(s_j)\) vanish.
Indeed the centralizer of a matrix with distinct eigenvalues
is commutative and contains the block idempotent.
The remaining even centralizer has dimension \(n\).
A fixed finite set of the defining linear equations has these
same kernels over \(F\).
Their Smith factors bound finite-precision kernels by
\(Q^C\) times \(Q^{m}\) for each surviving generic dimension,
which gives exactly (B1r).
The Frobenius approximation in the tower transports these
bounds to the actual new local groups for large \(N\).

\paragraph{The finite duality count and the exponent cancellation.}
Let \(\operatorname{Sel}_{1,N,m}\) and
\(\operatorname{Bd}_{1,N,m}\) be the actual primal cocycle and
boundary sets.
We give the count with every term visible.
At an ordinary place take local cocycles equal to a boundary
plus an element of \(Z_{b,v,m}\).
The map from a correcting matrix and an upper cocycle onto
this sum has fibers containing the upper-triangular module,
of cardinality \(Q^{m(n+\nu)}\).
Thus, writing \(M_m=W_1/T^m\),
\[
       \#\mathcal L^Z_v
               \leq \#M_m\,
                       \#Z_{b,v,m}/Q^{m(n+\nu)} .
 \tag{B1s}
\]
At an old non-\(\ell\) place the unramified cocycle set has
cardinality \(\#M_m\), since inertia acts trivially.
At the distinguished and the new places take the full local
cocycle set.
Local duality at a non-\(\ell\) place gives its cardinality
as \(\#M_m\) times the cardinality of the dual invariants.
These are bounded by (B1r).

Finally each real place has at least \(Q^{m\nu}\) invariants
in \(M_m\).
For the chosen complex conjugation these are its skew matrices,
of rank \(\nu\).
For another real conjugation, the polarized representation
identifies its fixed space with that skew space by multiplication
by the corresponding symmetric invertible matrix.
This works integrally, so the same lower bound holds at
finite precision.

Apply the finite global-duality cardinal formula, using the
perfect pairing in (B1d), the local subgroups just specified,
and (B1o).
In cocycle notation the local factors are
\(\#\mathcal L^Z_w/\#M_m\); the real factors occur in the
denominator, and the global primal invariants in the numerator.
The denominator of global dual invariants can be discarded
for an upper bound.
Normalization at complex conjugation identifies the primal
class quotient with
\(\operatorname{Sel}_{1,N,m}/\operatorname{Bd}_{1,N,m}\),
because \(2\) is invertible.
Using (B1q)--(B1s) therefore gives
\[
 \begin{split}
 \#\operatorname{Sel}_{1,N,m}\,
   Q^{m[(n+\nu)\#\{v\mid\ell\}+\nu d]}
 \ \leq\ 
 \#\operatorname{Bd}_{1,N,m}\,
 Q^{C_G+C_{\rm tr}+C_*+\sum_vC_v+\sum_jC_j}\\
 {}\cdot
 Q^{\,m\sum_{v\mid\ell}
        [n(d'_v+1)+\nu(d_v+1)]+mnr}.
 \end{split}
 \tag{B1t}
\]
All cohomology and cocycle sets here are finite by the
arithmetic finiteness input.
The auxiliary level \(N\) only needs to exceed a threshold
depending on \(m\).

The identity (B1p) gives the exact cancellation
\[
 \sum_v[n(d'_v+1)+\nu(d_v+1)]+nr
       -[(n+\nu)\#\{v\mid\ell\}+\nu d]=nd+nr .
\]
Hence one constant \(C\), fixed before every precision and
tower level, satisfies
\[
 \boxed{\ 
 \#\operatorname{Sel}_{1,N,m}
 \leq\#\operatorname{Bd}_{1,N,m}\,
                       Q^{\,m(nd+nr)+C}
 \quad\text{for all sufficiently large \(N\), for each \(m\)}.
       \ }
 \tag{B1u}
\]
This is \srcid{TWGrp.B1} with its actual parameter
\(D_0^{\rm bound}=nd+nr\).
The point-dependent bounded constants affect only the
precision threshold and additive constant; the displayed
linear slope is explicit.

\paragraph{Exact sources and what this closes.}
All paths in this paragraph are under
\srcid{Development.AutomorphicArithmetic.AutomorphicVarieties}
at \srcid{alpoge-source-7cbf7a0c}.
The odd relation calculation is
\srcid{PatchingModules.DeterminantPadicFree.RelPair.relPair},
\srcid{RelMove.relMove}, \srcid{DTw.der},
\srcid{DerEv.derEv}, \srcid{DTwAlt.eigen},
\srcid{Infres.infres}, \srcid{Shapiro.shapiro},
and \srcid{BlockPrinc.ann}; these are assembled in
\srcid{OddPart.oddPart}, \srcid{EvenPart.evenPart},
and \srcid{DetAlg.detAlg}.
Integral denominator clearing is \srcid{DetAsm.detAsm}.
The actual dense and finite choices are
\srcid{AutomorphicDetection.Detection.detection},
\srcid{DeterminantPadicFree.MemPowSmul.z1Fg},
\srcid{AcConst.matFg}, \srcid{FinSub.finSub},
\srcid{DetFin.detFin}, \srcid{OmDense.omDense},
and \srcid{AutomorphicDetection.Choice.choice}.
Finite-precision detection is
\srcid{EpsilonEtaMatrices.ExistsNcardAnnihilatedCocyclesLeCoboundariesMul};
descent and normalization are
\srcid{Descend} and \srcid{DualN.dualN}.
The local and invariant estimates are
\srcid{EntryZero}, \srcid{Rho1MapReductionLevel},
\srcid{FinrankCentralizerW1Le}, and the local ordinary-count
modules imported by \srcid{OneLevel}.
The finite duality count is \srcid{OneLevel.oneLevel};
its arithmetic assembly is
\srcid{W1BoundArithmeticRankTwo.W1_bound_arithmetic_rank_two}
and \srcid{AutomorphicDetection.BoundW1Arithmetic.b1Ar}.
These sources are in
\url{https://github.com/ChristianKudeba/alpoge-rh-verification/tree/alpoge-source-7cbf7a0c/src/Development/AutomorphicArithmetic/AutomorphicVarieties}.

Thus the first- and second-Selmer estimates and the arithmetic
Taylor--Wiles group tower have explicit proofs once the actual
nice points, their regular-coset property, and the established
arithmetic cohomology are in place.
This result supplies those numerical and group-theoretic
producers to the completed-support construction.
It does not identify an arbitrary polarized representation
with an automorphic point; that realization remains the
separate classical and seed/support part of the argument.
